\subsection{Dimension of a ring, height of an ideal}

DEFINITION 5. — The Krull dimension, or simply the dimension, of a (commutative) ring A is the Krull dimension of the topological space , and is denoted by  (II, § 4, no. 3, definition 4). If  is a prime ideal of A, the dimension of A at \_ is denoted by , namely \_.

The map \(\mathfrak{p} \mapsto \mathrm{V}(\mathfrak{p})\) is a decreasing bijection from the set of prime ideals of A onto the set of irreducible closed subsets of \(\operatorname{Spec}(A)\) (loc. cit., cor. 2 to prop. 14). The dimension of A is therefore the supremum of the set of lengths of chains of prime ideals of A; it is equal to \(-\infty\) , \(+\infty\) or to a positive integer.

Let \(\mathfrak{p} \in \operatorname{Spec}(\mathbf{A})\); the sets \(\operatorname{Spec}(\mathbf{A})_f\), as \(f\) ranges over A, form a basis for the topology of \(\operatorname{Spec}(\mathbf{A})\), and \(\mathfrak{p}\) belongs to the open set \(\operatorname{Spec}(\mathbf{A})_f\) if and only if \(f\) does not belong to \(\mathfrak{p}\). Consequently, \(\dim_{\mathfrak{p}}(\mathbf{A})\) is the infimum of the numbers \(\dim(\mathbf{A}_f)\), where \(f\) ranges over \(\mathfrak{p}\) (II, § 5, no. 1, Proposition 1).

Examples. --- 1) We have  if and only if A is reduced to 0. For \textless to hold, it is necessary and sufficient that every prime ideal of A be maximal. The integral domains of dimension 0 are the fields. A Noetherian ring has dimension {} if and only if it is Artinian (IV, § 2, no. 5, prop. 9).

\begin{enumerate}
\def\labelenumi{\arabic{enumi})}
\setcounter{enumi}{1}
\item
  Dedekind rings are the Noetherian integrally closed rings of dimension \(\leqslant 1\) (VII, § 2, no. 2, th. 1). More generally, by V, § 1, no. 2, cor. 2 to prop. 9, a ring is a finite product of Dedekind rings if and only if it is Noetherian, reduced, integrally closed in its total ring of fractions, and of dimension \(\leqslant 1\) .
\item
  If A is a valuation ring (VI, § 1, No. 2, Def. 2), its dimension is equal to the height of the valuation (VI, § 4, No. 4, Prop. 5).
\item
  Let A be a ring. We have
\end{enumerate}

\[
\dim (\mathrm{A} [ \mathrm{X} ]) \geqslant \dim (\mathrm{A}) + 1.
\]

Indeed, if \(\mathfrak{p}_0 \subset \ldots \subset \mathfrak{p}_n\) is a chain of prime ideals of A, of length \(n\) , we obtain a chain \(\mathfrak{p}_0' \subset \ldots \subset \mathfrak{p}_{n+1}'\){}{} of prime ideals of A[X], of length \(n+1\) , by setting \(\mathfrak{p}_i' = \mathfrak{p}_iA[X]\) for \(0 \leqslant i \leqslant n\) , and \(\mathfrak{p}_{n+1}' = \mathfrak{p}_nA[X] + XA[X]\) .

By the same reasoning, one proves the inequality \(\dim(\mathrm{A}[[\mathrm{X}]]) \geqslant \dim(\mathrm{A}) + 1\) . From this one deduces by induction the inequalities

\[
\begin{array}{l} \dim (\mathrm{A} [ \mathrm{X} _ {1},..., \mathrm{X} _ {n} ]) \geqslant \dim (\mathrm{A}) + n, \\ \dim (\mathrm{A} [ [ \mathrm{X} _ {1},..., \mathrm{X} _ {n} ] ]) \geqslant \dim (\mathrm{A}) + n. \end{array}
\]

We shall show later (§ 3, no. 4, cor. 3 of prop. 7 and cor. 3 of prop. 8) that equality holds in the two preceding formulas when A is Noetherian.

* 5) Let X be a complex analytic manifold. If X has complex dimension n at a point x of X, the local ring of germs at x of analytic functions on X has dimension n.*

\begin{enumerate}
\def\labelenumi{\arabic{enumi})}
\setcounter{enumi}{5}
\item
  Let \(k\) be a field and A a nonzero integral \(k\)-algebra. Then one has \(\dim(A) = 0\) . This follows from cor. 1 of prop. 1 of V, § 2, no. 1, and from the fact that \(\dim(k) = 0\) .
\item
  If n is a nilideal of A, Spec(A) is homeomorphic to Spec(A/n) (II, § 4, no. 3, remark). Hence one has dim(A/n) = dim(A); in particular, one has dim(A) = dim(A\(_{red}\)) where A\(_{red}\) is the quotient of the ring A by its nilradical.
\item
  There exist Noetherian rings of infinite dimension (p. 83, exerc. 13). We shall see below (§ 3, no. 1, cor. 1 of prop. 2) that every Noetherian local ring has finite dimension.
\end{enumerate}

PROPOSITION 6. --- Let A be a ring.

\begin{enumerate}
\def\labelenumi{\alph{enumi})}
\item
  If \(\mathfrak{a}\) is an ideal of A, one has \(\dim (\mathrm{A} / \mathfrak{a})\leqslant \dim (\mathrm{A}).\)
\item
  If S is a multiplicative subset of A, one has \(\dim (\mathbf{S}^{-1}\mathbf{A})\leqslant \dim (\mathbf{A}).\)
\item
  We have \(\dim(A) = \sup_{\mathfrak{p}} \dim(A/\mathfrak{p})\) , where \(\mathfrak{p}\) runs through the set of minimal prime ideals of A.
\item
  If A has only finitely many minimal prime ideals (for example if A is Noetherian (II, § 4, no. 3, Corollary 3 to Proposition 14)) and if \(\mathfrak{p}\) is a prime ideal of A, one has
\end{enumerate}

\[
\dim_ {\mathfrak {p}} (\mathrm{A}) = \sup _ {\mathfrak {q}} \dim_ {\mathfrak {p} / \mathfrak {q}} (\mathrm{A} / \mathfrak {q}),
\]

where \(\mathfrak{q}\) runs through the set of minimal prime ideals of A contained in \(\mathfrak{p}\).

\begin{enumerate}
\def\labelenumi{\alph{enumi})}
\setcounter{enumi}{4}
\tightlist
\item
  Let  be an ideal of A which is contained in no minimal prime ideal of A; one then has . In particular, if A is an integral domain, one has  for every nonzero ideal  of A.
\end{enumerate}

By the remark in II, § 4, no. 3, if  is an ideal of A, the topological space  is homeomorphic to the closed subspace  of . Assertion  follows from this and from Proposition 1, a), of no. 1. Assertion  follows from loc. cit., the corollary to Proposition 13. By loc. cit., Corollary 2 to Proposition 14, the irreducible components of  are homeomorphic to the spaces , where  is a minimal prime ideal of A, and assertion  follows from the corollary of Proposition 1 of no. 1. Under the hypothesis of , the space  has only finitely many irreducible components; the irreducible components of  containing  are the sets , where  is a minimal prime ideal contained in . Assertion  then follows from the corollary, b), of Proposition 1 of no. 1.

Let us finally prove \(e\) ). We need to prove that, for any chain \(\mathfrak{p}_{0} \subset \ldots \subset \mathfrak{p}_{n}\) of prime ideals of A such that \(\mathfrak{a} \subset \mathfrak{p}_{0}\) , one has \(\dim(A) \geqslant n + 1\) . In view of the hypothesis made on \(\mathfrak{a}\) there exists a prime ideal \(\mathfrak{p}_{-1}\) of A contained in \(\mathfrak{p}_{0}\) , distinct from \(\mathfrak{p}_{0}\) , and \(\mathfrak{p}_{-1} \subset \mathfrak{p}_{0} \subset \ldots \subset \mathfrak{p}_{n}\) is a chain of prime ideals of A, of length \(n + 1\) .

Remark 1. --- Let \(\rho: A \to B\) be a ring homomorphism. Then  is the supremum of the numbers , where \(\rho(p)\) ranges over the set of minimal prime ideals of A: indeed, for every minimal prime ideal \(p\) of B, there exists a minimal prime ideal \(q\) of A contained in \(\mathfrak{p}\)\(\rho^{-1}(\mathfrak{q})\) (II, § 2, no. 6, lemma 2), and one has

\[
\dim (\mathbf {B} / \mathfrak {q}) \leqslant \dim (\mathbf {B} / \rho (\mathfrak {p})\mathbf {B}) \leqslant \dim (\mathbf {B})
\]

By prop. 6, a); we conclude by prop. 6, c).

DEFINITION 6. --- Let  be an ideal of a ring A. The codimension of  in  is called the height of the ideal  and is denoted by .

Suppose A is integral. Then the prime ideals of height 1 of A in the sense of def. 4 of VII, § 1, no. 6 are precisely the prime ideals of height 1 in the sense of the definition above.

PROPOSITION 7. --- a) The height of a prime ideal  of A is the supremum of the lengths of chains of prime ideals \(\mathfrak{p}_{0} \subset \ldots \subset \mathfrak{p}_{n}\) such that \(\mathfrak{p}_{n} = \mathfrak{p}\).

\begin{enumerate}
\def\labelenumi{\alph{enumi})}
\setcounter{enumi}{1}
\item
  Let  be a prime ideal of A and  an ideal of A. Then we have  if  is not contained in \(\dim(A_{p}/\mathfrak{a}A_{p}) = -\infty\) and  if  is contained in \(\dim(A_{p}/\mathfrak{a}A_{p}) = \text{codim}(V(p), V(a))\) . In particular, if  is a prime ideal of A, one has \(\dim(A_{\mathfrak{p}}) = \operatorname{ht}(\mathfrak{p})\) .
\item
  If  is an ideal of A, one has , where  ranges over the set of prime ideals of A containing \_.
\end{enumerate}

The assertion  is the translation of Definition 3, , of no. 2. The assertion  follows from the fact that the map  is an order-preserving isomorphism of the set of prime ideals  of A such that  onto the set of prime ideals of the local ring  (II, § 2, no. 5, Proposition 11). Let  be an ideal of A; the irreducible closed subsets of  are the sets , where  is a prime ideal of A containing . Assertion  therefore follows from Remark 1 of no. 2.

\(a\)\(a\)\(b\)\(\mathfrak{q} \mapsto \mathfrak{q}(A_{\mathfrak{p}} / aA_{\mathfrak{p}})\)\(\mathfrak{q}\)\(\mathfrak{a} \subset \mathfrak{q} \subset \mathfrak{p}\)\(A_{\mathfrak{p}} / aA_{\mathfrak{p}}\)\(\mathfrak{a}\)\(\mathfrak{a}\)COROLLARY. --- Let \(\mathfrak{p}\) be a prime ideal of A and S a multiplicative subset of A not meeting \(\mathfrak{p}\). Then \(\mathfrak{a}\)\(c\)

\(^{-1}\).

This follows from prop. 7, a), and II, § 2, no. 5, prop. 11.

PROPOSITION 8. --- Let A be a ring.

\begin{enumerate}
\def\labelenumi{\alph{enumi})}
\item
  We have  , where \(\dim(A) = \sup_{\mathfrak{m}} \dim(A_{\mathfrak{m}}) = \sup_{\mathfrak{m}} \operatorname{ht}(\mathfrak{m})\)\(m\)\(A\) runs through the set of maximal (resp. prime) ideals of A.
\item
  Let  be an ideal of A distinct from A and let  be an ideal of A contained in . Then we have  . In particular, for every ideal  of A distinct from A, we have the inequality .
\end{enumerate}

The first assertion follows from remark 2 of no. 2 and from prop. 7, b). The second follows from prop. 3, a) of no. 2 and from the relations , .

DEFINITION 7. --- A ring A is said to be catenary if the topological space Spec(A) is catenary (no. 2, def. 4).

This means therefore that, for every pair \((\mathfrak{p}, \mathfrak{q})\) of prime ideals of A such that \(\mathfrak{q} \subset \mathfrak{p}\) , all saturated chains of prime ideals of A with endpoints \(\mathfrak{p}\) and \(\mathfrak{q}\) have length \(\operatorname{codim}(\mathrm{V}(\mathfrak{p}), \mathrm{V}(\mathfrak{q})) = \dim(\mathrm{A}_{\mathfrak{p}} / \mathfrak{q}\mathrm{A}_{\mathfrak{p}})\) .

Remarks. --- 2) Every quotient ring of a catenary ring is catenary. For the ring A to be catenary, it is necessary and sufficient that, for every prime ideal  of A, the ring \(\mathfrak{p}\) be catenary.

\begin{enumerate}
\def\labelenumi{\arabic{enumi})}
\setcounter{enumi}{2}
\item
  By prop. 7, b) and prop. 4 of no. 2, the ring A is catenary if and only if, for every triple  of prime ideals of A such that , we have \(\dim(A_{p}/qA_{p}) + \dim(A_{q}/rA_{q}) = \dim(A_{p}/rA_{p})\) . If A is integral and of finite dimension, then A is catenary if and only if we have \(ht(q) + \dim(A_{p}/qA_{p}) = ht(p)\) for every pair  of prime ideals of A such that \((p, q)\)\(\mathfrak{q} \subset \mathfrak{p}\) . Indeed, the topological space Spec(A) is then irreducible and of finite dimension, and we apply the corollary to prop. 4 of no. 2.
\item
  Let A be a ring of finite dimension, all of whose maximal chains of prime ideals have the same length. Then A is catenary, we have  for every prime ideal  of A, and  for every pair \_ of prime ideals of A such that \_ (no. 2, prop. 5).
\item
  We shall see in § 2, no. 4, that every algebra of finite type over a field is a catenary ring. There exist local noetherian rings which are not catenary (p. 83, exerc. 16).
\end{enumerate}

